\section{PHP Project Analysis} % {{{

\marginnote{\includegraphics[width=\marginparwidth]{php/logo}}

PHP is a server side, interpreted programming language designed for web
applications and development \cite{PHPIntro}. As such, it requires a web server
with a connected PHP installation, though it can also be used with a command
line interface. It is widely available for all major platforms and is licensed
under the PHP License \cite{PHPManual}. It is worth noticing, that while it is
a Free Software license, it is incompatible to the widely used \ac{GPL}. PHP
comes with an extensive standard library and supports many programming
paradigms, such as object oriented, imperative or procedural programming
styles. The abbreviation PHP originally stood for Personal Home Page
\cite{PHPHistory}. With PHP version 3, the project changed it's name to PHP
Hypertext Processor. PHP is a very popular programming language for web
development and featuring a large number of websites
\cite{PHPW3Techs,PHPStats}.

\subsection{History} % {{{

In 1994, Rasmus Lerdorf created the first version of PHP for his own website,
naming the program Personal Home Page Tools \cite{PHPHistory}. He published the
bundle one year later, after rewriting the original version. That also included
the first name change to FI (Forms Interpreter) and then to Personal Home Page
Construction Kit the same year. At that time, PHP could be considered as a
advanced programming interface.

However, it received another makeover in 1996, combining the two previous names
to PHP/FI. This time, PHP truly began to evolve to a full featured programming
language. It included several modules for database or browser interaction. This
version was also known as PHP/FI 2.0. Following a popularity boom in 1997 and
1998 it reached it's limitations due to the design of the project and by being
almost solely developed by Lerdorf.

In 1997 Andi Gutmans and Zeev Suraski rewrote the PHP interpreter and
approached Lerdorf discussing the problems PHP had at that time. Together, they
rewrote and redesigned the language and also changed the name to PHP Hypertext
Processor. Version 3.0 of PHP was then released in 1998 and replaced PHP/FI 2.0
completely. On of the most important features of the new approach was to
provide an interface for modules, which can be used directly from the
programming language.

The design of the language was improved further in 1998 improving performance
and the modularity of modules and code. The new engine was called Zend Engine
and provided the core of PHP 4.0, which was released in 2000.

Four years later, PHP 5 was released featuring the successor Zend Engine 2.0.
This release was incompatible to the previous 4.0 version and a large
initiative was planned to promote the transition from PHP 4 to PHP 5.

% }}}

\subsection{Community} % {{{

The PHP project consists of a large user and developer community
\cite{Magnusson2010}. Additionally to the PHP interpreter, there are many
people working on PHP extensions and components, which get distributed
alongside the standard interpreter.

The PHP community provides a large amount of conferences and workshops to meet
in person and work together on PHP. While the \emph{International PHP
Conference} was the first of it's kind in 2001 \cite{PHPConferences}, there are
many opportunities to meet other PHP contributors. A lot of local PHP User
Groups exist and provide weekly or monthly meetings and meet-ups such as
workshops.

Most of the communication inside the PHP project is done through mailing lists
\cite{Magnusson2010}. It has mailing lists for most aspects of the project,
however the most important is the \emph{internals} mailing list. Most of the
development process, future ideas and new contributions are handled over that
list. Of course, there are also several \ac{IRC} channels and blogs of PHP
developers available.

\begin{figure}[htbp]
  \centering
  \includegraphics[width=0.8\textwidth]{php/commits_by_author}
  \caption{Monthly activity of the most active PHP Developers}
\end{figure}

Officially, there is no categorizing of developers and everyone is treated
equally. However the project uses a concept named karma, which basically means
that one increases his karma by contributing to the project with code,
discussions and new ideas \cite{Magnusson2010}. Furthermore, several people are
responsible for different parts of the code or modules and can be classified
using those criteria \cite{PHPCredits}.

\begin{figure}[htbp]
  \centering
  \includegraphics[width=0.8\textwidth]{php/commits_by_year}
  \caption{Yearly overview of commits to PHP}
\end{figure}

\paragraph{Developer}

One is able to call himself a PHP Developer, once he gets a PHP Subversion
account. According to the PHP community, an account needs to be earned, which
basically means that one has to provide several patches and contributions to
the PHP project first. Once one has shown enough commitment, he can apply for
an account. The account however will be only granted if a reference person
approves.

\paragraph{Module Author}

Each developer has the possibility to become a module author. This role has the
responsibility for a certain PHP module, which gets distributed with the PHP
release. Examples for such modules are database or imaging modules. For each
module, there can be more than one module authors, who share the
responsibility.

\paragraph{PHP Author}

This role defines a certain section of the PHP interpreter and it's
responsibility for it. It is quite similar to the module author role, except
that it covers only PHP itself.

\paragraph{Language Designer}

Currently only Rasmus Lerdorf, Andi Gutmans and Zeev Suraski hold this role.
They are the only who can actively change the programming language syntax and
concepts.

\begin{figure}[hbtp]
  \centering
  \includegraphics[width=\textwidth]{php/punchcard}
  \caption{Time based view on commits of contributors}
\end{figure}

% }}}

\subsection{Release Process} % {{{

In 2010, the PHP project set up a new release plan with detailed information
about how the release process should work \cite{PHPRelease}. Before,
individuals decided when a release happened and which features it included. The
new release cycle features two Release Managers, who are voted by PHP
Developers and are managing the release process. The PHP project uses a scheme
with three numbers for releases: major, minor and micro. Major releases get
published quite seldom, can break backwards compatibility and are planned a
long time in advance. Minor releases can have new features and bug fixes,
however backwards compatibility must be kept. Micro release are bugfix only
releases where backwards and API compatibility must be kept. Each major release
is followed by a number of minor releases, which themselves are followed by a
number of micro releases.

\begin{figure}[htbp]
  \centering
  \includegraphics[width=\textwidth]{php/releases}
  \caption{Major Releases of PHP}
\end{figure}

Starting with PHP 5.4 a major or minor release will get published each year. Each
yearly release is then followed by a number of micro releases for 2 years,
which contains bug fixes only. After that period, only micro releases with
security relevant fixes will get published.

\begin{figure}[htbp]
  \centering
  \includegraphics[width=0.8\textwidth]{php/release_plan}
  \caption{Preliminary PHP release cycle}
\end{figure}

% }}}

\subsection{Development} % {{{

Most of the development process is handled over the PHP mailing lists
\cite{PHPRelease,Magnusson2010,PHPVoting}. Not every patch however needs to be
discussed and approved first, minor features or changes often get directly
committed to the repository. Though every change gets through a peer review
process, where other developers review the made changes. Every change gets
reviewed by other developers, who in case the need arises, discuss the change
on the developers mailing list. They then often result in a more comprehensive
change or in new features.

\begin{figure}[htbp]
  \centering
  \includegraphics[width=0.8\textwidth]{php/commits_by_month}
  \caption{Amount of commits per month of Core Contributors}
\end{figure}

To help the development and decision process, the PHP project keeps track of
new ideas, features and proposals through so called \ac{RFC} \cite{PHPRFC}.
Those are documents which describe the new feature and it's rationale. A
\ac{RFC} document is always owned by at least one person, who is responsible
for it. The \ac{RFC} process begins with the author's submission of the
\ac{RFC} to the PHP wiki and an announcement to the internals mailing list.
Depending on the importance and affected sections of the PHP project, the
following discussion period is set to a minimum of one to two weeks. It can be
longer however, but not shorter \cite{PHPVoting}.

After the discussion period has passed, the author can either call for a vote
or extend the discussion period as needed \cite{PHPVoting}. The vote is
announced on the mailing list and followed by a voting period which should be
at least one week, but can be extended if required.

Depending on the importance of a \ac{RFC}, it needs either 2/3 of all votes or
\unit[50]{\%} plus one to get accepted. The importance is defined by whether it
actually changes the syntax or behaviour of the PHP language and is therefore
irreversible or not. A failed proposal can be resurrected at least six months
after the last vote or if the author makes considerably large changes to the
\ac{RFC}. There is no definition by the PHP project, what considerably large
changes constitute, however the opinion is that it should be changed in a
matter, that it significantly can change a vote's outcome.

\begin{figure}[htbp]
  \centering
  \begin{tikzpicture}[auto]
   \node[rectangle, rounded corners, draw=black, text width=6.5em, minimum height=2em, text centered]
     (draft) {Draft};
   \node[below=0.5cm of draft, rectangle, rounded corners, text width=6.5em, minimum height=2em, text centered]
     (dummy) {};
   \node[below=0.5cm of dummy, rectangle, rounded corners, draw=black, text width=6.5em, minimum height=2em, text centered]
     (withdrawn) {Withdrawn};

   \node[right=1cm of draft, rectangle, rounded corners, draw=black, text width=6.5em, minimum height=2em, text centered]
     (discussion) {Discussion};
   \node[below=0.5cm of discussion, rectangle, rounded corners, draw=black, text width=6.5em, minimum height=2em, text centered]
     (accepted) {Accepted};
   \node[below=0.5cm of accepted, rectangle, rounded corners, draw=black, text width=6.5em, minimum height=2em, text centered]
     (implemented) {Implemented};

   \node[right=1cm of discussion, rectangle, rounded corners, draw=black, text width=6.5em, minimum height=2em, text centered]
     (voting) {Voting};
   \node[below=0.5cm of voting, rectangle, rounded corners, draw=black, text width=6.5em, minimum height=2em, text centered]
     (declined) {Declined};

   \path[draw, -latex] (draft) -- (discussion);
   \path[draw, -latex] (discussion) -- (voting);
   \path[draw, -latex] ($(discussion.south) - (0.5cm,0)$) |-+(0,-0.2cm)-| (withdrawn);
   \path[draw, -latex] ($(voting.south) - (0.5cm,0)$) |-+(0,-0.2cm)-| (accepted);
   \path[draw, -latex] ($(voting.south) + (0.5cm,0)$) -- ($(declined.north) + (0.5cm,0)$);
   \path[draw, -latex] (accepted) -- (implemented);
   \path[draw, -latex] (declined.south) -- ++(0,-0.4cm) -- 
   ([yshift=-0.4cm, xshift=0.4cm] declined.south east) --
   ([yshift=0.4cm, xshift=0.4cm] voting.north east) -| (draft.north);
  \end{tikzpicture}
  \caption{Possible paths of the status of \ac{RFC} documents}
\end{figure}

All PHP developers are allowed to vote and if needed, representatives from
the PHP community, such as participants of PHP related discussions or
developers of PHP appendant projects can participate in the
voting process. Those however have to be chosen by PHP developers
\cite{PHPWhoVote}.

During the whole development process, a \ac{RFC} has a status assigned, which can be
one of the following \cite{PHPRFC}.

\paragraph{Draft}

Once a \ac{RFC} gets submitted to the PHP wiki, it gets this status assigned.
The Draft status means, that the \ac{RFC} is not ready yet for a discussion,
however it can be improved by other people as well and a preliminary discussion
can be started.

\paragraph{Discussion}

Once the author of a \ac{RFC} announces it for discussion, it gets this status
assigned. As previously described, the discussion phase starts with this
status.

\paragraph{Voting}

After the discussion phase the voting phase begins. While developers can vote
for or against the \ac{RFC}, the Voting status gets assigned including a link
to the voting site.

\paragraph{Accepted}

Once a voting was successful, the \ac{RFC} is accepted and will be implemented
or taken in action in near future.

\paragraph{Declined}

If the majority voted against the \ac{RFC}, it gets declined. After six months
and several changes to it, the author can resubmit it.

\paragraph{Implemented}

If the voting was successful, the \ac{RFC} will get implemented or taken in
action if it describes a community process.

\paragraph{Withdrawn}

The author can withdraw the \ac{RFC} without calling for a vote if he thinks
the \ac{RFC} is irrelevant or won't pass a vote.

\begin{figure}[htbp]
  \centering
  \includegraphics[width=0.8\textwidth]{php/authors_by_month}
  \caption{Amount of distinct authors of PHP over time}
\end{figure}

% }}}

% }}}
